% !TEX program = xelatex
\documentclass[12pt,a4paper]{article}

% --- Chinese support ---
\usepackage{xeCJK}
\setCJKmainfont{Songti SC}
\usepackage[margin=2.5cm]{geometry}
\usepackage{enumitem}
\usepackage{hyperref}
\usepackage{xcolor}
\usepackage{titlesec}

\titleformat{\section}{\Large\bfseries}{}{0em}{}
\titlespacing{\section}{0pt}{1em}{0.5em}
\setlength{\parskip}{4pt}

\title{答辩逐字稿 --- Defense Script\\
\large Asymptotic speed of the $N$-branching random walk with infinite offspring\\
\large 可无限分支的 $N$-分支随机游走的渐近速率性质}
\author{王亦扬 (Wang Yiyang)}
\date{2026年5月20日}

\begin{document}

\maketitle

% ==========================================
\section*{0. Opening / 开场}
% ==========================================

\begin{center}
\color{blue!60!black}\textbf{[Slide: Title Page]}
\end{center}

\begin{quote}
\textbf{各位老师早上好！我是王亦扬，我的硕士论文题目是\emph{可无限分支的 $N$-分支随机游走的渐近速率性质}，指导老师是 \'{E}lie A\"id\'ekon 教授。今天我将向大家汇报论文的主要结果和证明思路。}

\bigskip

\textbf{Good morning, distinguished committee members. I am Wang Yiyang. My master's thesis is titled \emph{Asymptotic speed of the $N$-branching random walk with infinite offspring}, advised by Professor \'{E}lie A\"id\'ekon. Today I will present the main results and proof ideas.}
\end{quote}

\medskip

\begin{center}
\color{blue!60!black}\textbf{[Slide: Outline]}
\end{center}

\begin{quote}
\textbf{报告分为六个部分：首先是模型背景与主要定理，然后是上界证明、下界证明，最后是最技术性的 $\mathrm{L}^2$ 极限定理与右尾/左尾估计。}

\bigskip

\textbf{The presentation has six parts: first, background and main theorem; then the upper bound proof; the lower bound; and finally the $\mathrm{L}^2$ limit theorem, which is the most technical part, together with right- and left-tail estimates.}
\end{quote}

% ==========================================
\section*{1. Background and Model / 模型背景}
% ==========================================

\begin{center}
\color{blue!60!black}\textbf{[Slide: $N$-BRW on the marked tree space --- Definition]}
\end{center}

\begin{quote}
\textbf{$N$-分支随机游走（$N$-BRW）的演化规则：从 $\lfloor N^a\rfloor$ 个粒子出发（$a\in[0,1]$）；每代粒子按一般分支律 $\Xi$ 独立繁殖，$\Xi$ 可以有可数无限个后代；从所有子代中仅保留位置最靠左的 $N$ 个粒子。关键困难在于 $\Xi$ 是无穷后代的一般分布，而非以往文献中的特定 Poisson 模型。左边的 3-BRW 示意图展示了选择机制。}

\bigskip

\textbf{The $N$-branching random walk ($N$-BRW) evolves by: start with $\lfloor N^a\rfloor$ particles at $0$ ($a\in[0,1]$); each particle reproduces according to a point process $\Xi$, which may have countably infinite offspring; keep only the $N$ leftmost particles among all offspring. The key difficulty is that $\Xi$ is a general distribution with possibly infinite offspring, not the specific Poisson model of earlier works. The left diagram is a 3-BRW illustrating selection.}
\end{quote}

\begin{center}
\color{blue!60!black}\textbf{[Slide: $N$-BRW on the marked tree space --- Formalism]}
\end{center}

\begin{quote}
\textbf{用有标记树（Ulam--Harris tree）严格表述：$\mathbb V$ 是顶点空间，$T\subset\mathbb V$ 是有根标号树，$(T,V)$ 是标记树给每个顶点赋予实数位置。分支随机游走 $\mathcal V$ 是随机标记树。选择算子 $\Phi_N$ 是确定性泛函：输入一群标记树，保留每代最靠左的 $N$ 个粒子。关键量为前沿位置 $\mathcal M_n^{N|a}$ 和渐近速度 $v_N$。速度的存在性由 Kingman 次加性遍历定理保证。}

\bigskip

\textbf{Formally: $\mathbb V$ is the Ulam--Harris vertex space; $T\subset\mathbb V$ is a rooted labelled tree; $(T,V)$ is a marked tree assigning real positions to vertices. A BRW is a random marked tree $\mathcal V$. The selection operator $\Phi_N$ is a deterministic functional: it keeps the $N$ leftmost particles per generation. Key quantities: the leftmost position $\mathcal M_n^{N|a}$ and the asymptotic speed $v_N$. Existence of $v_N$ follows from Kingman's subadditive ergodic theorem.}
\end{quote}

\begin{center}
\color{blue!60!black}\textbf{[Slide: Brunet--Derrida Conjecture]}
\end{center}

\begin{quote}
\textbf{Brunet 和 Derrida (1997) 发现：有限种群选择导致扩张速度减慢，速度修正为 $v_\infty-v_N\asymp(\log N)^{-2}$。这个 $(\log N)^{-2}$ 标度律出现在极值统计、对数相关场、平均场自旋玻璃等多个领域，具有普适性。本文将该标度律推广到一般分支律。}

\bigskip

\textbf{Brunet and Derrida (1997) discovered that finite-population selection causes a slowdown: $v_\infty-v_N\asymp(\log N)^{-2}$. This $(\log N)^{-2}$ scaling appears in extreme value statistics, log-correlated fields, and mean-field spin glasses---it is universal. My thesis extends this law to general branching laws.}
\end{quote}

% ==========================================
\section*{2. Main Theorem / 主要定理}
% ==========================================

\begin{center}
\color{blue!60!black}\textbf{[Slides: Main Theorem, Assumptions, Boundary Case, Existence]}
\end{center}

\begin{quote}
\textbf{主要结果：在边界情形（$\psi(1)=\psi'(1)=0$）下，渐近速度为}
\[
v_N = \frac{\pi^2}{2}\sigma^2\frac{1}{\log^2 N} + o_N\!\left(\frac{1}{\log^2 N}\right),
\]
\textbf{其中 $\sigma^2=\psi''(1)$ 是脊柱步长方差。系数 $\frac{\pi^2}{2}\sigma^2$ 是显式的、普适的，仅依赖于方差。边界情形意味着 size-biased 脊柱行走是中心化的，此时无约束 BRW 的速度为零（$v_\infty=0$），加性鞅退化，需要用导数鞅 $D_n$ 来重整化。Jaffuel (2009) 表明一般情况下 BRW 可变换到边界情形，故边界情形的假设不损失一般性。速度 $v_N$ 的存在性由 Kingman 定理保证。}

\bigskip

\textbf{Main result: in the boundary case ($\psi(1)=\psi'(1)=0$), the asymptotic speed is}
\[
v_N = \frac{\pi^2}{2}\sigma^2\frac{1}{\log^2 N} + o_N\!\left(\frac{1}{\log^2 N}\right),
\]
\textbf{where $\sigma^2=\psi''(1)$ is the spine-step variance. The coefficient $\frac{\pi^2}{2}\sigma^2$ is explicit and universal, depending only on the variance. The boundary case means the size-biased spine walk is centered; then $v_\infty=0$, the additive martingale degenerates, and renormalization requires the derivative martingale $D_n$. By Jaffuel (2009), a BRW can generally be transformed to the boundary case, so the assumption is without loss of generality. Existence of $v_N$ follows from Kingman's theorem.}
\end{quote}

% ==========================================
\section*{3. Upper Bound / 上界证明}
% ==========================================

\begin{center}
\color{blue!60!black}\textbf{[Slide: Upper Comparison Construction]}
\end{center}

\begin{quote}
\textbf{先证上界：$v_N\le\frac{\pi^2}{2}\sigma^2(\log N)^{-2}$。构造一个``比较过程''：杀死除最左侧外的所有粒子，让该粒子独自演化 $\ell_N$ 代作为 $N$-BRW，保留最左侧粒子，重复。这个过程显然比原始 $N$-BRW 更靠右。定义偏序 $\succcurlyeq$：A 控制 B 当且仅当 A 的粒子数更少且整体位置更靠右。重启过程在偏序下控制原始过程，给出 $\mathcal M_{k\ell_N}^{N|a}\le\mathcal M_{k\ell_N}^{\mathrm{upper}}$。}

\bigskip

\textbf{For the upper bound: $v_N\le\frac{\pi^2}{2}\sigma^2(\log N)^{-2}$. Construct a comparison process: kill all but the leftmost particle, let it evolve as an $N$-BRW of depth $\ell_N$, keep the leftmost, and repeat. This process is obviously to the right of the original $N$-BRW. Define partial order $\succcurlyeq$: A dominates B if A has fewer particles and is shifted right. The restart process dominates the original in this order, giving $\mathcal M_{k\ell_N}^{N|a}\le\mathcal M_{k\ell_N}^{\mathrm{upper}}$.}
\end{quote}

\begin{center}
\color{blue!60!black}\textbf{[Slides: SLLN, What Needs to Be Proved, Restart, Decomposition, Bad Event, Conclusion]}
\end{center}

\begin{quote}
\textbf{由强大数律，每个周期增量为 i.i.d.，得 $v_N\le\frac{1}{\ell_N}\E[\mathcal M_{\ell_N}^{N|0}]$。取 $\ell_N=t\log^3 N$，问题归结为证明 $\E[\mathcal M_{\ell_N}^{N|0}]\le\frac{\pi^2}{2}\sigma^2 t\log N+o(\log N)$。这里用到了 Theorem 1.1（$\mathrm{L}^2$ 极限定理）。但 Theorem 1.1 对 $a=0$ 需要指数矩条件，我们引入\emph{重启机制}来规避：将 $\Xi$ 截断为至多二分叉过程，用备份粒子提供 $\lfloor N^\alpha\rfloor$ 个后代，将 $a=0$ 转化为 $a=\alpha>0$，从而去除指数矩假设。上界证明分为三步：(1) 定义好事，将期望分解为三项；(2) 前两项由 Proposition 3.1 和 Wald 恒等式控制；(3) 坏事概率由分支小偏差估计证明 $\P(E_N^{\mathrm{c}})\le N^{-\varepsilon}$。最终对 $t$ 取 $\inf$ 得 $(\log N)^{-2}$ 上界。}

\bigskip

\textbf{By SLLN, the cycle increments are i.i.d., yielding $v_N\le\frac{1}{\ell_N}\E[\mathcal M_{\ell_N}^{N|0}]$. With $\ell_N=t\log^3 N$, the task reduces to $\E[\mathcal M_{\ell_N}^{N|0}]\le\frac{\pi^2}{2}\sigma^2 t\log N+o(\log N)$. This invokes Theorem 1.1 (the $\mathrm{L}^2$ limit theorem). But Theorem 1.1 for $a=0$ requires an exponential moment. We introduce a \emph{restart mechanism} to bypass this: truncate $\Xi$ to an at-most binary process, use backup particles to supply $\lfloor N^\alpha\rfloor$ offspring, reducing $a=0$ to $a=\alpha>0$, thus removing the exponential moment condition. The upper bound proceeds in three steps: (1) define a good event and decompose the expectation into three terms; (2) bound the first two by Proposition 3.1 and Wald's identity; (3) bound the bad event probability via small deviations: $\P(E_N^{\mathrm{c}})\le N^{-\varepsilon}$. Taking $\inf_t$ yields the $(\log N)^{-2}$ upper bound.}
\end{quote}

% ==========================================
\section*{4. Lower Bound / 下界证明}
% ==========================================

\begin{center}
\color{blue!60!black}\textbf{[Slides: Lower Comparison, SLLN]}
\end{center}

\begin{quote}
\textbf{下界由镜像构造证明：取 $\widetilde N = \lfloor N - N/\log N\rfloor$ 构造一个稍小的 $\widetilde N$-BRW，将其嵌入原始 $N$-BRW，应用相同的 SLLN 论证配上 Theorem 1.1，得到匹配的下界。上下界夹逼给出精确的渐近速度公式。}

\bigskip

\textbf{The lower bound uses a mirror construction: take $\widetilde N = \lfloor N - N/\log N\rfloor$, construct a slightly smaller $\widetilde N$-BRW, embed it into the original $N$-BRW, and apply the same SLLN argument with Theorem 1.1 to obtain a matching lower bound. The sandwich yields the exact asymptotic speed formula.}
\end{quote}

% ==========================================
\section*{5. $\mathrm{L}^2$ Limit Theorem / $\mathrm{L}^2$ 极限定理}
% ==========================================

\begin{center}
\color{blue!60!black}\textbf{[Slides: Theorem 1.1, Killed BRW Construction, Coupling, Right Tail Steps 1--2, Cases $a=1,a=0$]}
\end{center}

\begin{quote}
\textbf{Theorem 1.1 是技术核心：在 $\ell_N/\log^3 N\to t$ 的标度下，}
\[
\frac{1}{\log N}\min_{|u|=\ell_N}\mathcal V^{N|a}(u)\;\xrightarrow{\mathrm{L}^2}\;-a+\frac{\pi^2}{2}\sigma^2 t.
\]
\textbf{$\mathrm{L}^2$ 收敛归结为两个尾估计：$\P(\mathcal M^N_{\ell_N}\ge \zeta'\log N)\le N^{-\varepsilon}$ 和 $\P(\mathcal M^N_{\ell_N}\le \zeta\log N)\le N^{-\varepsilon}$。右尾的处理：不用选择而用``杀死''——构造被杀 BRW，设 Phase I（管道 $[0,\ell]$，宽 $b\log N$）和 Phase II（管道 $[\ell,\ell_N]$，宽 $cL^{1/3}$），粒子离开管道即被杀。参数取 $\ell=\log^{5/2}N$，$L=\ell_N-\ell$。通过 Lemma 3.17 将被杀 BRW 与 $N$-BRW 耦合，将右尾分解为 $\P(\tau\le\ell_N)+\P(\mathcal M^{\mathrm{k}}_{\ell_N}\ge \zeta'\log N)$。Step 1 用小偏差控制灭绝时间 $\tau$。Step 2 用 Lemma 3.9 和 Lemma 3.10 分别处理管道终点``粒子多''和``粒子少''两种情形。对于 $a=1$ 直接约化，$a=0$ 用重启机制。}

\bigskip

\textbf{Theorem 1.1 is the technical core: under $\ell_N/\log^3 N\to t$,}
\[
\frac{1}{\log N}\min_{|u|=\ell_N}\mathcal V^{N|a}(u)\;\xrightarrow{\mathrm{L}^2}\;-a+\frac{\pi^2}{2}\sigma^2 t.
\]
\textbf{$\mathrm{L}^2$ convergence reduces to two tail bounds: $\P(\mathcal M^N_{\ell_N}\ge \zeta'\log N)\le N^{-\varepsilon}$ and $\P(\mathcal M^N_{\ell_N}\le \zeta\log N)\le N^{-\varepsilon}$. For the right tail, we replace selection by killing: construct a killed BRW with Phase I (tube $[0,\ell]$, width $b\log N$) and Phase II (tube $[\ell,\ell_N]$, width $cL^{1/3}$); particles are killed upon leaving the tube. Parameters: $\ell=\log^{5/2}N$, $L=\ell_N-\ell$. Lemma 3.17 couples the killed BRW to the $N$-BRW, decomposing the right tail as $\P(\tau\le\ell_N)+\P(\mathcal M^{\mathrm{k}}_{\ell_N}\ge \zeta'\log N)$. Step 1 controls extinction time $\tau$ via small deviations. Step 2 handles the ``many particles'' and ``few particles'' cases at the tube endpoint via Lemma 3.9 and Lemma 3.10. The case $a=1$ reduces directly; $a=0$ uses the restart mechanism.}
\end{quote}

\begin{center}
\color{blue!60!black}\textbf{[Slides: Key Lemmas, Many-to-One, Spine, Left Tail, Mogul'skii]}
\end{center}

\begin{quote}
\textbf{核心引理串：Many-to-One Lemma 将分支期望转为脊柱期望——在 size-biased 测度下，沿脊柱位移 i.i.d.，其余粒子条件分布同原始 BRW。Lemmas 3.4--3.7 给出水平和倾斜管道的小偏差界。Lemma 3.8 是二阶矩界，Assumption 6 在此直接使用，并通过 Paley--Zygmund 传播到 Lemma 3.9 和 Lemma 3.10。左尾估计使用``大跳跃''事件：一个粒子向左异常跳跃，破坏管道约束，得到多项式衰减 $N^{-\varepsilon}$。}

\bigskip

\textbf{Key lemmas: The Many-to-One Lemma converts branching expectations to spine expectations---under the size-biased measure, spine increments are i.i.d., and other particles are conditionally distributed as the original BRW. Lemmas 3.4--3.7 give small deviation bounds in horizontal and tilted tubes. Lemma 3.8 is a second-moment bound; Assumption 6 is used directly here and propagates to Lemma 3.9 and Lemma 3.10 via Paley--Zygmund. The left tail uses a ``large-drop'' event: one particle makes an unusually large leftward jump, breaking the tube constraint, yielding polynomial decay $N^{-\varepsilon}$.}
\end{quote}

\begin{quote}
\textbf{最后讲 $\frac{\pi^2}{2}$ 从何而来。Mogul'skii 小偏差定理：中心化随机游走在宽 $\Delta$ 的管道中保持 $L$ 步的概率 $\sim\exp(-\frac{\pi^2}{2}\sigma^2\frac{L}{\Delta^2})$（$\Delta^2\ll L$）。这与 Cram\'er 大偏差完全不同——大偏差给指数衰减，小偏差给亚指数修正。$\frac{\pi^2}{2}$ 是 Dirichlet 拉普拉斯算子在 $[0,1]$ 上的第一特征值，是一个谱量。本题取 $\Delta=cL^{1/3}$，代入 Mogul'skii 后经参数优化，$\frac{\pi^2}{2}\sigma^2$ 自然成为主导系数。}

\bigskip

\textbf{Finally, where does $\frac{\pi^2}{2}$ come from? The Mogul'skii small deviation theorem: the probability that a centered random walk stays in a tube of width $\Delta$ for $L$ steps satisfies $\sim\exp(-\frac{\pi^2}{2}\sigma^2\frac{L}{\Delta^2})$ ($\Delta^2\ll L$). This differs fundamentally from Cram\'er large deviation---the latter gives exponential decay, while small deviation gives a sub-exponential correction. $\frac{\pi^2}{2}$ is the principal eigenvalue of the Dirichlet Laplacian on $[0,1]$, a spectral quantity. Taking $\Delta=cL^{1/3}$ and substituting into Mogul'skii yields $\frac{\pi^2}{2}\sigma^2$ as the leading coefficient after parameter optimization.}
\end{quote}

% ==========================================
\section*{6. Summary / 总结}
% ==========================================

\begin{center}
\color{blue!60!black}\textbf{[Slide: Thank You]}
\end{center}

\begin{quote}
\textbf{总结本文三个主要贡献：(1) 将 A\"id\'ekon--Hu 的 $\mathrm{L}^2$ 极限定理从特定 Poisson 分支律推广到一般分支律（可数无限后代）；(2) 通过重启机制去除了 $a=0$ 情形的指数矩条件，构造了截断+重启策略；(3) 严格验证了 $(\log N)^{-2}$ 标度律，给出包含显式常数 $\frac{\pi^2}{2}\sigma^2$ 的渐近速度公式，确认了该标度律在一般分支律下的普适性。谢谢各位老师，请各位老师批评指正。}

\bigskip

\textbf{To summarize, three main contributions: (1) Extended the $\mathrm{L}^2$ limit theorem of A\"id\'ekon--Hu from specific Poisson branching to general branching laws with possibly infinite offspring; (2) Removed the exponential moment condition for $a=0$ via a restart mechanism with binary truncation; (3) Rigorously verified the $(\log N)^{-2}$ scaling law with an explicit constant $\frac{\pi^2}{2}\sigma^2$, confirming universality under general branching laws. Thank you for your attention. I am happy to answer any questions.}
\end{quote}

\end{document}
